\documentclass[
bibliography=totoc,
headings=big,
captions=tableheading,
chapterprefix=true% like in standard class "report"
]{scrreprt}
\usepackage[utf8]{inputenc}
\usepackage[frenchb]{babel}
\usepackage[left=2cm,right=2cm,top=1.5cm,bottom=3cm]{geometry}
\usepackage{xcolor}
\usepackage{graphicx}
\usepackage{minitoc}
\usepackage{setspace}
\usepackage{multicol}

% package for symbol-fonts
\usepackage{marvosym}

% other fonts
\usepackage{mathpazo}

% Finally, give us PDF bookmarks
\usepackage{color,hyperref}
\definecolor{darkblue}{rgb}{0.0,0.0,0.3}
\hypersetup{colorlinks,breaklinks,
            linkcolor=darkblue,urlcolor=darkblue,
            anchorcolor=darkblue,citecolor=darkblue}

\begin{document}
    \pagenumbering{gobble}

\hfill \includegraphics[scale=0.1]{goelug}

\section*{Présentation de l'association GoéLUG}

% \dictum[Sun Tzu]{If you know the enemy and you know yourself %\\
% you need not fear the results of a hundred battles.}

\minitoc

GoéLUG est la combinaison de goéland, oiseau très connu des havrais, et LUG, acronyme anglo-saxon qui désigne un groupe d’utilisateurs Linux. 
C'est une association loi 1901 à fonctionnement collégial dont les buts affichés sont:
    
\begin{itemize}
\item   promouvoir l'utilisation des logiciels libres sous toutes ses formes : débats, installs parties, cours ou ateliers.

 \item   démocratiser les technologies de l’Internet, ce dernier fonctionnant essentiellement grâce au logiciel libre.

\item    échanger autour de la culture libre et de l’art libre.

\end{itemize}


\section*{Notre besoin}

Un lieu pour organiser des événements de manière régulière et/ou ponctuelle, comme :
    

\begin{enumerate}

\item des install parties (10 à 30 personnes). 

    Le principe d'une install party est d'accompagnement des personnes pour l'installation d'un système d'exploitation libre, type Ubuntu. 

\item des permanences (5 à 10 personnes)        

\item des ateliers (5 à 10 personnes): 

\begin{enumerate}

\item    GoéBidouille: ateliers autour de hardware et d'électronique (exemple: raspberry)

\item    Initiation au développement de logiciels et de sites web

\item    Initiation logiciels libres (traitement de texte, mails, graphisme, ...)

\end{enumerate}

\end{enumerate}
        

\section*{Ce que nous pourrions offrir pour la cantine numérique 'Le Container'}

    % Assistance aux associations dans leur stratégie de communication (montage d’un site Web, hébergement,…) <= je ne suis pas certain de cette proposition. / Ok, je vais la enlever pour le lettre

\begin{itemize}

\item    Assistance à la demande, de déploiement de serveurs, d’installation de systèmes d’exploitation libres, \dots

\item  Atelier pédagogique et vulgarisation du monde libre en informatique ( réunion, install Party, Crypto party, \dots) en soirée ou en week-end.

\item Sécuriser votre informatique dans le World Wild Web.

\end{itemize}

\section*{Un calendrier}

\begin{multicols}{2}

\subsubsection*{Permanence:}

Proposition de calendrier

\begin{itemize}

\item jeudi 7 janvier 2016

\item jeudi 4 février 2016

\item jeudi 3 mars 2016 

\item jeudi 7 avril 2016 

\item jeudi 5 mai 2016

\item jeudi 2 juin 2016

\end{itemize}

Proposition d'horaires : 18h-22h.
% Ce qui correspondes a une formule facile a retenir comme "chaque premiere jeudi du mois".

\columnbreak

\subsubsection*{Install parties:}

Ponctuelles, environ 4 fois par an, si possible le samedi toute la journée.

\subsubsection*{Ateliers:}

Ponctuels, si possible le samedi à la demi-journée.

\end{multicols}
\hfill \href{http://www.goelug.org/}{\tt www.goelug.org} / \href{mailto:contact@goelug.org}{\tt contact@goelug.org}

\end{document}
